\chapter{Processos, Troca de Contexto e Escalonamento}

\chapterauthor{João Pedro Rizzo Sales, Samuel Romano Ribeiro, Guilherme Michelsen do Monte}


\section{Descrição}


Este capítulo documenta três módulos do UFSKernel:
o módulo de \textbf{processos}, responsável por representar, criar e encerrar
unidades de execução; o mecanismo de \textbf{troca de contexto}, que realiza
a alternância segura entre modo usuário e modo kernel; e o \textbf{escalonador},
que decide qual processo recebe a CPU e por quanto tempo. Juntos, esses módulos
viabilizam a multiprogramação no kernel, conforme descrito em
Silberschatz et al.~\cite{silberschatz2018} e Tanenbaum e Bos~\cite{tanenbaum2015}.

Os arquivos que compõem esses módulos são: \texttt{process.h},
\texttt{process.c}, \texttt{fork.c}, \texttt{exec.c}, \texttt{exit.c},
\texttt{proc\_ll.c}, \texttt{svc\_entry.S}, \texttt{irq\_entry.S},
\texttt{irq\_return.S}, \texttt{context\_switch.S}, \texttt{scheduler.h} e
\texttt{process\_control\_block.c}.


\section{Fundamentação Teórica}


\subsection{Processo e PCB}

Um processo é um programa em execução: uma entidade dinâmica que engloba
o contador de programa, os registradores, a pilha e o espaço de endereços
ocupado na memória~\cite{tanenbaum2015}. Para gerenciá-los, o kernel mantém
um \textbf{Bloco de Controle de Processo} (PCB) por processo, concentrando
tudo o que é necessário para suspendê-lo e retomá-lo como se nunca tivesse
sido interrompido~\cite{silberschatz2018}.

Durante seu ciclo de vida, um processo transita entre quatro estados. Os três
canônicos, \textbf{Pronto}, \textbf{Em execução} e \textbf{Bloqueado}, são
descritos em~\cite{tanenbaum2015}. O UFSKernel acrescenta um quarto,
\textbf{Zumbi}: o processo encerrou, mas sua estrutura ainda não foi
descartada, pois o processo pai deve registrar o encerramento do filho antes
que os recursos sejam liberados~\cite{silberschatz2018}.

\subsection{Trapframe, Contexto de Kernel e ABI ARM}

Na ARMv7, o processador opera em modos de privilégio distintos. Toda transição
de modo usuário para kernel, seja por chamada de sistema (\texttt{SVC}) ou por
interrupção (\texttt{IRQ}), exige salvar o estado completo do processo para que
a execução possa ser retomada. A estrutura que armazena
esse estado é o \textbf{trapframe}~\cite{tanenbaum2015}.

Para trocas de contexto internas ao kernel, quando o escalonador substitui
um processo por outro sem sair do modo supervisor, apenas um subconjunto
dos registradores precisa ser salvo. A ABI ARM (AAPCS)~\cite{arm2014} define
dois grupos: registradores \emph{call-preserved} (r4--r11, SP e LR), que uma
função deve restaurar antes de retornar, e registradores \emph{call-clobbered}
(r0--r3, r12), que o chamador sabe que podem ser destruídos. Como o compilador
já protege os call-clobbered em torno de qualquer chamada de função, o
mecanismo de troca de contexto de kernel precisa salvar apenas os call-preserved.

\subsection{Escalonamento: Round-Robin e MLFQ}

O escalonador decide qual processo pronto recebe a CPU e por quanto tempo
(o \textbf{quantum}), tipicamente acionado por um timer periódico~\cite{silberschatz2018}.
O \textbf{Round-Robin} organiza os processos em fila circular com quantum fixo,
mas trata igualmente processos interativos e CPU-bound~\cite{tanenbaum2015}.
A \textbf{Multilevel Feedback Queue} (MLFQ) supera essa limitação: mantém
múltiplas filas de prioridade com quanta crescentes, promovendo processos que
usam pouco CPU (interativos) e rebaixando os que esgotam o quantum (CPU-bound),
de modo que cada processo migra naturalmente para o nível compatível com seu
comportamento real~\cite{silberschatz2018}.


\section{Módulo de Processos}


\subsection{Decisões de Arquitetura e Design}

O módulo mantém duas estruturas distintas para o estado de execução de um
processo. O \texttt{struct trapframe} armazena todos os registradores de
propósito geral (r0--r12), os registradores banked do modo usuário
(\texttt{sp\_usr}, \texttt{lr\_usr}), o contador de programa (\texttt{pc\_usr})
e o registrador de estado (\texttt{cpsr\_usr}) — tudo o que o processo de
usuário enxerga. O \texttt{struct context} armazena apenas os registradores
call-preserved (r4--r11, SP e LR), suficientes para a troca de contexto
dentro do kernel. Salvar r0--r3 e r12 no contexto de kernel seria redundante,
pois o compilador já os trata como voláteis em torno de qualquer
chamada~\cite{arm2014}.

\subsubsection{Alocação estática de memória}

Sem um alocador dinâmico disponível, toda a memória dos processos foi
provisionada em tempo de compilação com três arrays globais em \texttt{fork.c}:

\lstset{captionpos=t}
\begin{lstlisting}[language=C, caption={Arrays de alocação estática em \texttt{fork.c}}]
struct process structProcessos[MAX_PROCESS_COUNT];   // 64 PCBs
char memProcessos[MAX_PROCESS_COUNT][SIZE_1MB];      // 64 x 1 MB
char kstackProcessos[MAX_PROCESS_COUNT][SIZE_16KB];  // 64 x 16 KB
\end{lstlisting}
\lstset{captionpos=b}

Isso impõe limites fixos: máximo de 64 processos simultâneos, 1~MB de
memória e 16~KB de pilha de kernel por processo, independentemente da
necessidade.

\subsubsection{Simulação de executáveis com \texttt{struct arquivo}}

Sem um sistema de arquivos disponível, a noção de executável foi simulada
pela \texttt{struct arquivo}: um ponteiro de função (\texttt{start}) e um
tamanho (\texttt{size}).

\lstset{captionpos=t}
\begin{lstlisting}[language=C, caption={Estrutura de simulação de arquivo em \texttt{process.h}}]
struct arquivo {
    void (*start)(void); // funcao que representa o programa
    int size;            // tamanho do codigo a copiar
};
\end{lstlisting}
\lstset{captionpos=b}

Isso permite testar o ciclo completo fork--exec--exit sem depender de
sistema de arquivos, mas implica que o código do processo coexiste na mesma
imagem binária do kernel, sem separação real de espaço de endereços.

\subsubsection{Gestão de PIDs e a estrutura \texttt{process\_block\_list}}

PIDs são gerenciados por \texttt{create\_pid} e \texttt{free\_pid} sobre
o array circular \texttt{used\_pids[MAX\_PROCESS\_COUNT]}. O PID~0 é
reservado, convenção Unix que evita confundi-lo com um ponteiro
nulo~\cite{silberschatz2018}. A estrutura \texttt{process\_block\_list},
definida em \texttt{process.h}, prevê filas de processos bloqueados por
\emph{sleep} ou por operações de E/S, mas não é utilizada na implementação
atual, servindo como ponto de extensão planejado~\cite{tanenbaum2015}.

\subsection{Detalhes da Implementação}

\subsubsection{O PCB: \texttt{struct process}}

\lstset{captionpos=t}
\begin{lstlisting}[language=C, caption={Bloco de Controle de Processo em \texttt{process.h}}]
struct process {
    char* mem;              // inicio da regiao de memoria do processo
    unsigned msize;         // tamanho da regiao de memoria
    char* kstack;           // inicio da pilha de kernel
    struct trapframe *tf;   // ponteiro para o trapframe (offset 12)
    enum State state;       // RUNNING, READY, BLOCKED ou ZOMBIE
    enum Block_Reason blocked_by;
    pid_t pid;
    process parent;
    pid_list children_ids;
    struct context context; // registradores call-preserved (troca de contexto)
    unsigned priority;
};
\end{lstlisting}
\lstset{captionpos=b}

A ordem dos quatro primeiros campos é deliberada: \texttt{mem} (4~bytes),
\texttt{msize} (4~bytes) e \texttt{kstack} (4~bytes) posicionam \texttt{tf}
no offset~12 — valor codificado na constante \texttt{TRAPFRAME\_OFFSET = 12}
usada pelo código assembly para acessá-lo diretamente.

\subsubsection{Ciclo de vida: \texttt{first\_process}, \texttt{fork}, \texttt{exec} e \texttt{exit}}

\texttt{first\_process} cria o processo inicial: aloca a estrutura por
indexação estática no PID gerado, insere o processo no escalonador, define
\texttt{current}, posiciona o trapframe no topo da kstack e chama
\texttt{exec}. Ao final, chama \texttt{irq\_return} para entregar o controle
ao modo usuário como se fosse um retorno de interrupção.

\texttt{sys\_fork} copia byte a byte a memória e a kstack do pai para o
filho, recalcula a posição do trapframe e altera os ponteiros de contexto
do usuário (\texttt{pc\_usr}, \texttt{lr\_usr}, \texttt{sp\_usr}) para o
espaço de endereços do filho:

\lstset{captionpos=t}
\begin{lstlisting}[language=C, caption={Translação de endereços no \texttt{sys\_fork}}]
unsigned int pc_offset =
    newprocess->parent->tf->pc_usr
    - (unsigned int)newprocess->parent->mem;
newprocess->tf->pc_usr =
    (unsigned int)newprocess->mem + pc_offset;
\end{lstlisting}
\lstset{captionpos=b}

O pai recebe o PID do filho em \texttt{r0}; o filho recebe zero, convenção
Unix que permite distinguir os dois após o retorno~\cite{silberschatz2018}.
O \texttt{context.lr} do filho aponta para \texttt{irq\_return}, garantindo
que o primeiro \texttt{context\_switch} para ele desembarque corretamente
no retorno ao modo usuário.

\texttt{sys\_exec} copia o conteúdo do pseudo-arquivo para a memória do
processo e prepara um trapframe novo. O endereço de entrada é calculado como:

\lstset{captionpos=t}
\begin{lstlisting}[language=C, caption={Cálculo do entry point em \texttt{sys\_exec}}]
unsigned int entry_point =
    (unsigned int) p->mem
    + ((((unsigned int)arquivo->start) & ~1U) - 0x40000000);
\end{lstlisting}
\lstset{captionpos=b}

A operação \texttt{\& \textasciitilde{}1U} zera o bit menos significativo
do ponteiro de função — necessário porque em Thumb esse bit é 1 para sinalizar
o modo de instrução, mas não faz parte do endereço real~\cite{arm2014}.
O \texttt{cpsr\_usr} é configurado para \texttt{0x30}: modo USR com Thumb ativo.

\texttt{sys\_exit} remove o processo do escalonador, libera o PID, marca o
estado como \texttt{ZOMBIE} e invoca \texttt{pcb\_elect} para eleger
imediatamente o próximo processo.


\section{Mecanismo de Troca de Contexto}


\subsection{Três mecanismos: entrada, retorno e troca}

O mecanismo de troca de contexto é composto por três partes complementares:
a \textbf{entrada} (\texttt{svc\_entry.S}, \texttt{irq\_entry.S}), que captura
o estado do processo de usuário e o armazena em um trapframe na pilha de kernel;
o \textbf{retorno} (\texttt{irq\_return.S}), que restaura esse estado e devolve
o controle ao processo; e a \textbf{troca de contexto de kernel}
(\texttt{context\_switch.S}), que salva e restaura apenas os registradores
call-preserved para alternar entre processos enquanto permanece em modo supervisor.
Mantê-los separados permite ao escalonador decidir, a cada interrupção, se
retorna ao mesmo processo ou realiza uma troca.

\subsection{Entrada via SVC e IRQ}

Ao receber um SVC ou IRQ, o hardware ARMv7 salva automaticamente o CPSR em
SPSR e o endereço de retorno em LR~\cite{arm2014}. O problema é que
r0--r12 ainda contêm valores do processo interrompido, mas é preciso um
registrador livre para operar. A solução é gravar r0 e r1 em posições
negativas da pilha IRQ antes de reutilizá-los:

\lstset{captionpos=t}
\begin{lstlisting}[language={[x86masm]Assembler},
                   caption={Salvamento temporário e troca de pilha em \texttt{irq\_entry.S}}]
str r0, [sp, #-4]            @ salva r0 temporariamente na pilha IRQ
str r1, [sp, #-8]            @ salva r1 temporariamente na pilha IRQ
mov r1, sp                   @ guarda o SP_irq para restaurar r0 e r1 depois

ldr r0, =current             @ carrega endereco da variavel global current
ldr r0, [r0]                 @ r0 = ponteiro para o processo atual
ldr sp, [r0, #TRAPFRAME_OFFSET] @ troca pilha IRQ pela kstack do processo

ldr r0, [r1, #-4]            @ restaura r0
ldr r1, [r1, #-8]            @ restaura r1
\end{lstlisting}
\lstset{captionpos=b}

Em seguida, o trapframe é construído na kstack do processo:

\lstset{captionpos=t}
\begin{lstlisting}[language={[x86masm]Assembler},
                   caption={Construção do trapframe na pilha de kernel}]
sub sp, sp, #(17 * 4)        @ reserva espaco para os 17 campos do trapframe
stmia sp, {r0-r12}           @ salva r0-r12
add r12, sp, #(13*4)
stmia r12, {sp, lr}^         @ salva sp_usr e lr_usr (^ acessa banked USR)
str lr, [sp, #(15 * 4)]      @ salva PC (LR contem o end. de retorno do HW)
mrs r12, spsr
str r12, [sp, #(16 * 4)]     @ salva CPSR do processo (estava em SPSR)
\end{lstlisting}
\lstset{captionpos=b}

O sufixo \texttt{\^} em \texttt{stmia r12, \{sp, lr\}\^} é essencial: ele
instrui o processador a acessar os registradores banked do modo usuário
(\texttt{SP\_usr} e \texttt{LR\_usr}), e não os do modo corrente. Sem ele,
o estado do processo seria corrompido~\cite{arm2014}.

\subsection{Retorno ao Modo Usuário}

\texttt{irq\_return} realiza o processo inverso: lê \texttt{tf} do processo
atual, restaura o CPSR do usuário no SPSR, carrega r0--r12, restaura
\texttt{sp\_usr} e \texttt{lr\_usr} com o sufixo \texttt{\^} e executa:

\lstset{captionpos=t}
\begin{lstlisting}[language={[x86masm]Assembler},
                   caption={Retorno ao modo usuário em \texttt{irq\_return.S}}]
movs pc, lr
\end{lstlisting}
\lstset{captionpos=b}

O sufixo \texttt{s} com o PC como destino em modo privilegiado copia
automaticamente o SPSR para o CPSR~\cite{arm2014} — restaurando modo,
flags e conjunto de instruções em uma única instrução.

\subsection{Troca de Contexto de Kernel}

\lstset{captionpos=t}
\begin{lstlisting}[language={[x86masm]Assembler},
                   caption={Troca de contexto de kernel em \texttt{context\_switch.S}}]
context_switch:
    cmp r0, #0
    beq if_null              @ se nao ha processo saindo, apenas carrega

    stmia r0!, {r4-r11}      @ salva registradores call-preserved
    str sp, [r0], #4         @ salva SP
    str lr, [r0]             @ salva LR (ponto de retorno quando reativado)

if_null:
    ldmia r1!, {r4-r11}      @ restaura registradores call-preserved
    ldr sp, [r1], #4         @ restaura SP (agora na pilha do novo processo)
    ldr lr, [r1]             @ restaura LR

    bx lr                    @ retorna -- no contexto do processo entrante
\end{lstlisting}
\lstset{captionpos=b}

O caso \texttt{r0 == NULL} cobre o boot, quando não há processo saindo.
A propriedade central é que \texttt{bx lr} não retorna para quem chamou
\texttt{context\_switch}, mas para o ponto em que o processo entrante havia
sido suspenso. Para um processo recém-criado via \texttt{fork}, esse ponto é
\texttt{irq\_return}, colocado ali por \texttt{sys\_fork}, fazendo-o
ingressar diretamente no retorno ao modo usuário.

\section{Escalonador}

\subsection{Decisões de Arquitetura e Design}

\subsubsection{MLFQ com 32 níveis e quanta crescentes}

A ideia adotada é uma MLFQ simplificada com 32 níveis. O \texttt{pcb}
é um array de \texttt{struct pcb\_entry} com quanta variando de 20 a 175
em progressão aritmética de passo 5. Processos recém-criados entram no
nível~0, maior prioridade, menor quantum, refletindo a heurística de que
processos novos tendem a ser interativos~\cite{tanenbaum2015}. A progressão
aritmética produz separação gradual entre os níveis sem saltos abruptos que
pudessem causar \emph{starvation} nos níveis inferiores.

\subsubsection{Margem de quantum e \texttt{list\_location}}

A constante \texttt{SCHEDULER\_QUANTUM\_MARGIN = 5} define uma zona de
tolerância: sem ela, pequenas variações no tempo acumulado causariam
oscilações de promoção e rebaixamento. A variável global \texttt{list\_location}, atualizada por
\texttt{pcb\_get\_node\_pid}, comunica o nível atual
do processo para \texttt{pcb\_climb} e \texttt{pcb\_fall}. É uma escolha que evita parâmetros extras, mas cria acoplamento implícito:
uma chamada intercalada a \texttt{pcb\_get\_node\_pid} com outro PID
sobrescreveria \texttt{list\_location} e corromperia a movimentação. Em um
sistema multiprocessado, isso seria uma fonte de \emph{race conditions}.

\subsubsection{Alocação estática e \texttt{time\_elapsed}}

Os nós da lista ligada são alocados estaticamente em
\texttt{process\_blocks\_nodes[64]}, indexados por PID, impondo o mesmo
limite de 64 processos. O comportamento do escalonador é
funcional, mas o feedback de prioridade não reflete tempo real.

\subsection{Detalhes da Implementação}

\subsubsection{Estruturas de dados}

\lstset{captionpos=t}
\begin{lstlisting}[language=C, caption={Estruturas do escalonador em \texttt{scheduler.h}}]
struct pcb_entry {
    pll_node *process_list; // cabeca da lista ligada de processos
    unsigned quantum;       // quantum deste nivel
    short process_count;    // numero de processos no nivel
    short next_process;     // indice circular do proximo a ser eleito
};

typedef struct process_node_ll {
    struct process_node_ll *next;
    process proc;
} pll_node;
\end{lstlisting}
\lstset{captionpos=b}

O campo \texttt{next\_process} implementa o revezamento circular dentro de
cada nível: após eleger o processo de índice $i$, o campo avança para
$(i+1) \bmod \texttt{process\_count}$.

\subsubsection{O algoritmo de eleição: \texttt{pcb\_elect}}

Invocada a cada \emph{tick} do timer via \texttt{timer\_callback},
\texttt{pcb\_elect} executa em três fases:

\textbf{Fase 1 — Acumular tempo.} Adiciona o retorno de \texttt{time\_elapsed}
à variável \texttt{running\_for}.

\textbf{Fase 2 — Avaliar o processo atual.} Se o processo ainda está dentro
do quantum (com a margem), retorna sem preemptar:

\lstset{captionpos=t}
\begin{lstlisting}[language=C, caption={Avaliação do quantum em \texttt{pcb\_elect}}]
if (current->state == RUNNING &&
    running_for < (quantum + SCHEDULER_QUANTUM_MARGIN))
    return; // ainda dentro do quantum, nao preempta

// esgotou o quantum: aplica feedback de prioridade
if (running_for <= quantum - SCHEDULER_QUANTUM_MARGIN)
    pcb_climb(current_process_node); // usou pouco: promove
else
    pcb_fall(current_process_node);  // usou muito: rebaixa
\end{lstlisting}
\lstset{captionpos=b}

\textbf{Fase 3 — Eleger o próximo.} Percorre os 32 níveis do mais alto para
o mais baixo. Em cada nível, localiza o processo na posição
\texttt{next\_process} e avança enquanto o processo não estiver
\texttt{READY}. Ao encontrar um candidato, atualiza \texttt{current},
marca-o como \texttt{RUNNING}, zera \texttt{running\_for} e chama
\texttt{context\_switch}. O \texttt{goto end} permite abandonar o laço
interno sem flags booleanas adicionais.

\subsubsection{Promoção e rebaixamento}

\texttt{pcb\_climb} e \texttt{pcb\_fall} movem o processo um nível acima
ou abaixo, respectivamente, removendo o nó do nível atual e inserindo-o no
adjacente. Os limites (nível~0 e nível~31) são respeitados: processos nos
extremos não são movidos.


\section{Interface Pública (API)}


\subsection{Módulo de Processos (\texttt{process.h})}

\begin{description}
    \item[\texttt{pid\_t create\_pid()}] Aloca e retorna um PID disponível.
    Retorna \texttt{MAX\_PROCESS\_COUNT} se todos estiverem em uso.

    \item[\texttt{void free\_pid(pid\_t pid)}] Libera o PID para reutilização.

    \item[\texttt{void first\_process(struct arquivo arquivo)}] Cria e lança
    o processo inicial. Deve ser chamada uma única vez pelo \texttt{kmain}.
    Não retorna.

    \item[\texttt{int fork()}] Cria um processo filho. Retorna o PID do filho
    ao pai e 0 ao filho (wrapper assembly em \texttt{syscall.S}).

    \item[\texttt{void exec(struct arquivo* arquivo)}] Substitui a imagem do
    processo atual. Não retorna ao chamador.

    \item[\texttt{void exit()}] Encerra o processo atual. Não retorna.
\end{description}

\subsection{Módulo do Escalonador (\texttt{scheduler.h})}

\begin{description}
    \item[\texttt{void pcb\_elect(void)}] Seleciona e ativa o próximo
    processo. Chamada automaticamente pelo timer; pode ser invocada
    explicitamente por \texttt{sys\_exit}.

    \item[\texttt{void pcb\_insert(unsigned priority\_level, pll\_node* proc)}]
    Insere um processo no nível especificado. Nível~0 é o de maior prioridade.

    \item[\texttt{void pcb\_remove(unsigned priority\_level, pll\_node* proc)}]
    Remove um processo do nível especificado.

    \item[\texttt{pll\_node* pcb\_get\_node\_pid(pid\_t pid)}] Localiza o nó
    do processo pelo PID. Efeito colateral: atualiza \texttt{list\_location}.

    \item[\texttt{void pcb\_climb(pll\_node* proc)}] Move o processo um nível
    acima. Requer chamada prévia de \texttt{pcb\_get\_node\_pid} para o
    mesmo processo.

    \item[\texttt{void pcb\_fall(pll\_node* proc)}] Move o processo um nível
    abaixo. Mesma precondição de \texttt{pcb\_climb}.

    \item[\texttt{pll\_node* pll\_node\_new(process proc)}] Inicializa e
    retorna o nó de lista ligada alocado estaticamente para o processo,
    indexado pelo PID.

    \item[\texttt{int proc\_is\_valid(process proc)}] Retorna verdadeiro se
    o processo não é nulo e não está em estado \texttt{ZOMBIE}.
\end{description}


%-------------------------------------------------------------------------------
\section{Testes e Verificação}
%-------------------------------------------------------------------------------

\subsection{Rotina de Teste}

O \texttt{kmain} chama \texttt{first\_process(arqinicio)}, que executa
\texttt{programainicio}:

\begin{enumerate}
    \item Imprime \texttt{"Processo rodando programa inicio"}.
    \item Chama \texttt{fork}; o filho invoca \texttt{exec(\&arq1)},
    carregando \texttt{programa1}.
    \item O pai chama \texttt{fork} novamente; o segundo filho invoca
    \texttt{exec(\&arq2)}, carregando \texttt{programa2}.
    \item O pai chama \texttt{exit}.
\end{enumerate}

Os dois filhos entram em laço infinito, imprimindo suas mensagens com
um delay de espera ativa. O escalonador alterna entre eles a cada
\emph{tick} do timer.

\subsection{Saída Esperada}

\lstset{captionpos=t}
\begin{lstlisting}[caption={Saída serial esperada no QEMU}]
Bem-vindo ao UFSKernel!
Executando em modo ARM bare-metal no QEMU.
Preparando CPU (VBAR e Pilha IRQ)...
Configurando GIC...
Ligando interrupcoes...
Ativando Timer...
Kernel rodando na main...
Processo rodando programa inicio
Processo rodando programa 1
Processo rodando programa 1
Processo rodando programa 1
Processo rodando programa 1
Processo rodando programa 1
>>> TICK DO TIMER <<<
Processo rodando programa 2
Processo rodando programa 2
Processo rodando programa 2
Processo rodando programa 2
Processo rodando programa 2
>>> TICK DO TIMER <<<
Processo rodando programa 1
...
\end{lstlisting}
\lstset{captionpos=b}

A intercalação de \texttt{">>> TICK DO TIMER <<<"} com as impressões dos
processos confirma que o escalonador está sendo acionado corretamente e
alternando entre os processos a cada quantum.
